\documentclass[12pt]{report}
\usepackage[utf8]{inputenc}
\usepackage[spanish]{babel}
\usepackage{geometry}
\geometry{margin=2.5cm}
\usepackage{graphicx}
\usepackage{hyperref}
\hypersetup{
    colorlinks=true,
    linkcolor=black,
    urlcolor=blue,
    citecolor=black
}
\usepackage{booktabs}
\usepackage{array}
\usepackage{longtable}
\usepackage{float}
\usepackage{enumitem}
\usepackage{amssymb}
\usepackage{titlesec}
\usepackage{fancyhdr}
\usepackage{caption}
\usepackage{subcaption}
\usepackage{amsmath}
\usepackage{listings}
\usepackage{xcolor}

% Estilo de código C
\lstdefinestyle{cstyle}{
    language=C,
    basicstyle=\ttfamily\small,
    keywordstyle=\color{blue}\bfseries,
    commentstyle=\color{gray}\itshape,
    stringstyle=\color{orange},
    numbers=left,
    numberstyle=\tiny\color{gray},
    breaklines=true,
    frame=single,
    backgroundcolor=\color{gray!5},
    tabsize=4
}

\pagestyle{fancy}
\fancyhf{}
\rhead{SmartRoad Rover}
\lhead{Electrónica Digital III}
\cfoot{\thepage}

\begin{document}

% ============================================================
% PORTADA
% ============================================================
\begin{titlepage}
    \centering
    \vspace*{1cm}
    \includegraphics[width=0.3\textwidth]{Escudo UdeA.jpg}\\[1.5cm]
    {\large \textbf{Electrónica Digital III}}\\[0.5cm]
    {\Huge \textbf{Avance del Proyecto Final}}\\[0.5cm]
    {\LARGE \textbf{SmartRoad Rover}}\\[1cm]
    Juan Felipe Orozco Londoño\\[0.3cm]
    Universidad de Antioquia\\[2cm]
    \today
\end{titlepage}

% ============================================================
% ÍNDICE
% ============================================================
\tableofcontents
\newpage

% ============================================================
% CAPÍTULO 1 — PLAN DE TRABAJO
% ============================================================
\chapter{Plan de Trabajo}

\section{Descripción general}

El proyecto SmartRoad Rover es desarrollado de forma individual por Juan Felipe Orozco Londoño. A continuación se presenta el plan de trabajo estructurado en 6 semanas, con las tareas definidas, su descripción y el estado de avance actual. El lenguaje de programación utilizado para el firmware es \textbf{C}, haciendo uso del \textbf{Pico SDK} oficial de Raspberry Pi, con una arquitectura basada en \textbf{polling e interrupciones (IRQ)}.

\section{Desglose de tareas y avance}

\subsection*{Semana 1 — Definición y diseño del sistema \hfill \textcolor{green!60!black}{100\%}}

\begin{itemize}
    \item Definición de requisitos funcionales y no funcionales.
    \item Selección y justificación de componentes.
    \item Elaboración del diagrama de bloques del hardware.
    \item Estimación del presupuesto.
    \item Configuración del repositorio en GitHub con estructura de carpetas.
\end{itemize}

\subsection*{Semana 2 — Diseño detallado de hardware y arquitectura de software \hfill \textcolor{green!60!black}{100\%}}

\begin{itemize}
    \item Diseño del esquemático del circuito (Pico W, HC-SR04, TB6612FNG, NEO-6M).
    \item Definición de conexiones y asignación de pines del microcontrolador.
    \item Diseño del prototipo sobre protoboard.
    \item Definición de la arquitectura de software en C/Pico SDK: módulos, máquina de estados, flujo de polling e interrupciones.
    \item Documentación del diseño en el repositorio.
\end{itemize}

\subsection*{Semana 3 — Adquisición de componentes y montaje del prototipo \hfill \textcolor{orange}{60\%}}

\begin{itemize}
    \item Adquisición de sensores HC-SR04, driver TB6612FNG, módulo GPS NEO-6M, chasis 2WD, batería 18650 y regulador step-down. \textit{(En proceso)}
    \item Montaje del circuito sobre protoboard. \textit{(Pendiente)}
    \item Verificación de alimentación y continuidad del circuito. \textit{(Pendiente)}
    \item Prueba básica de encendido del sistema. \textit{(Pendiente)}
\end{itemize}

\subsection*{Semana 4 — Desarrollo e integración del firmware \hfill \textcolor{orange!80!red}{15\%}}

\begin{itemize}
    \item Configuración del entorno de desarrollo Pico SDK en C.
    \item Implementación del módulo de lectura de sensores HC-SR04 mediante interrupciones (IRQ en flanco de subida y bajada).
    \item Implementación del módulo de control de motores DC vía PWM.
    \item Implementación del servidor web embebido para control manual vía WiFi.
    \item Implementación de la máquina de estados para gestión de modos de operación.
    \item Integración de módulos y pruebas unitarias.
\end{itemize}

\subsection*{Semana 5 — Integración completa y pruebas del sistema \hfill \textcolor{red}{0\%}}

\begin{itemize}
    \item Integración del firmware completo sobre el prototipo ensamblado.
    \item Pruebas del modo autónomo reactivo (detección, evasión, tiempos de respuesta verificados por UART).
    \item Pruebas del modo manual vía WiFi.
    \item Pruebas del modo GPS en entorno exterior.
    \item Verificación de comportamientos de contingencia (pérdida de WiFi, pérdida de señal GPS).
    \item Corrección de errores y ajustes finos.
\end{itemize}

\subsection*{Semana 6 — Validación final y preparación de la entrega \hfill \textcolor{red}{0\%}}

\begin{itemize}
    \item Pruebas completas en el escenario de pruebas definitivo (entorno exterior, 10 m $\times$ 10 m).
    \item Verificación de todos los requisitos funcionales (RF) y no funcionales (RNF).
    \item Actualización final del repositorio: firmware, hardware, documentación.
    \item Preparación de la sustentación del proyecto.
\end{itemize}

\section{Resumen de avance global}

\begin{table}[H]
\centering
\renewcommand{\arraystretch}{1.4}
\begin{tabular}{|l|c|}
\hline
\textbf{Fase} & \textbf{Avance} \\
\hline
Diseño y planificación del sistema & 100\% \\
Diseño de hardware y esquemático & 100\% \\
Adquisición de componentes y montaje & 60\% \\
Desarrollo de firmware (C / Pico SDK) & 15\% \\
Integración y pruebas & 0\% \\
Validación final y entrega & 0\% \\
\hline
\textbf{Avance general estimado} & \textbf{$\approx$45\%} \\
\hline
\end{tabular}
\caption{Resumen del estado de avance del proyecto}
\end{table}

% ============================================================
% CAPÍTULO 2 — CAMBIOS RESPECTO A LA RONDA 2
% ============================================================
\chapter{Cambios Respecto a la Definición — Ronda 2}

Se describen a continuación todos los cambios aplicados al proyecto respecto a la entrega de la Ronda 2, incluyendo correcciones derivadas de la retroalimentación del docente y decisiones técnicas tomadas durante la fase de diseño detallado.

\section{Lenguaje de programación: MicroPython $\rightarrow$ C / Pico SDK}

El firmware del sistema se implementará en \textbf{C utilizando el Pico SDK oficial} de Raspberry Pi. Esta decisión responde al requerimiento explícito del curso y es técnicamente más adecuada para una arquitectura basada en polling e interrupciones (IRQ), ya que el SDK en C ofrece acceso directo al hardware, control determinista de tiempos y manejo nativo de periféricos como PWM, UART y GPIO sin capas de abstracción intermedias.

\section{Implementación física: PCB $\rightarrow$ Protoboard}

La implementación física del circuito se realizará sobre \textbf{protoboard} en lugar de una placa de circuito impreso (PCB), por las siguientes razones:

\begin{itemize}
    \item El tiempo disponible hasta la entrega final no es suficiente para completar el ciclo de diseño de layout, fabricación, recepción y soldadura de una PCB con garantías de calidad.
    \item La protoboard permite iteraciones rápidas durante el desarrollo y la depuración del firmware, facilitando modificaciones de conexiones sin riesgo de dañar un circuito permanente.
    \item Para un prototipo funcional de validación académica, la protoboard cumple completamente con todos los requisitos funcionales y no funcionales definidos.
\end{itemize}

El esquemático y el presupuesto de fabricación de PCB se conservan como documentación técnica, representando la hoja de ruta para un escalamiento futuro del proyecto. Como complemento, se contempla el diseño de una \textbf{carcasa en impresión 3D} para alojar el conjunto sobre el chasis 2WD, aportando robustez mecánica al prototipo.

\section{Actualización del presupuesto}

Como consecuencia del cambio de PCB a protoboard, el presupuesto se actualiza de la siguiente manera:

\begin{table}[H]
\centering
\renewcommand{\arraystretch}{1.4}
\begin{tabular}{|l|c|}
\hline
\textbf{Categoría} & \textbf{Subtotal (COP)} \\
\hline
Componentes electrónicos & 126.000 \\
Protoboard y materiales de conexión & 25.000 \\
Plataforma mecánica & 53.000 \\
Sistema de alimentación & 48.000 \\
\hline
\textbf{Total estimado} & \textbf{252.000} \\
\hline
\end{tabular}
\caption{Presupuesto actualizado — prototipo sobre protoboard}
\end{table}

El total se reduce de 312.000 COP a \textbf{252.000 COP}, liberando recursos para la adquisición de componentes faltantes.

\section{Correcciones de la Ronda 2 incorporadas}

Los siguientes ajustes fueron aplicados con base en la retroalimentación del docente:

\begin{itemize}
    \item \textbf{Nota administrativa:} se confirmó que no hubo cambio de nombre del proyecto ni novedades en el equipo de trabajo.
    \item \textbf{RNF reformulados:} el tiempo de respuesta se separó por modo de operación. El modo autónomo reactivo tiene un ciclo determinista menor a 200 ms; el modo manual reconoce que la latencia WiFi no es controlable por el microcontrolador.
    \item \textbf{RNF de tolerancia a fallos agregados:} RNF-5 (detención en $\leq$ 500 ms ante pérdida de WiFi) y RNF-6 (transición a modo autónomo reactivo ante pérdida de señal GPS por más de 3 segundos).
    \item \textbf{Escenario de pruebas unificado:} único entorno exterior de mínimo 10$\times$10 m, materiales de obstáculos especificados (madera o cartón rígido) y verificación de tiempos por timestamps UART en lugar de cronómetro manual.
    \item \textbf{Radio GPS corregido:} reducido de 3 m a 1,5 m para ser coherente con el área de pruebas y la precisión real del módulo NEO-6M.
    \item \textbf{Costo de envío PCB incluido:} partida de envío internacional ($\sim$20.000 COP) incorporada al presupuesto de referencia para escalamiento futuro.
\end{itemize}

% ============================================================
% CAPÍTULO 3 — CIRCUITOS ANALÓGICOS
% ============================================================
\chapter{Circuitos Analógicos Requeridos}

El sistema SmartRoad Rover requiere \textbf{cuatro circuitos analógicos} de acondicionamiento. Tres de ellos son idénticos (un divisor de tensión por cada sensor HC-SR04) y uno corresponde al sistema de alimentación con regulación de voltaje. A continuación se presenta la lista, el procedimiento de diseño de cada uno y su estado de avance.

\section{Lista de circuitos requeridos}

\begin{enumerate}
    \item Divisor de tensión resistivo — sensor HC-SR04 frontal (ECHO: 5 V → 3.3 V)
    \item Divisor de tensión resistivo — sensor HC-SR04 izquierdo (ECHO: 5 V → 3.3 V)
    \item Divisor de tensión resistivo — sensor HC-SR04 derecho (ECHO: 5 V → 3.3 V)
    \item Circuito de alimentación — regulador step-down 7.4 V → 5 V con dos dominios de voltaje
\end{enumerate}

\section{Circuitos 1, 2 y 3 — Divisor de tensión resistivo (HC-SR04 ECHO)}

\subsection{Problema que resuelve}

El Raspberry Pi Pico W opera con lógica de \textbf{3.3 V}. Sus pines GPIO soportan un máximo de 3.3 V en la entrada. El sensor HC-SR04 opera a \textbf{5 V} y devuelve la señal ECHO a nivel de 5 V. Conectar directamente el pin ECHO al GPIO del Pico W daña el microcontrolador de forma permanente. Es necesario adaptar el nivel lógico de 5 V a 3.3 V antes de que la señal llegue al GPIO.

El pin TRIG del HC-SR04 acepta 3.3 V como nivel alto válido, por lo tanto la señal de salida del Pico W va directamente al TRIG sin circuito intermedio.

\subsection{Topología seleccionada}

Divisor resistivo de tensión. Es la solución más simple, económica y confiable para adaptación de nivel lógico unidireccional (el ECHO solo va del sensor al microcontrolador, no al revés). No requiere componentes activos ni encapsulados especiales.

\subsection{Procedimiento de diseño y cálculo}

La ecuación del divisor de tensión es:

\[
V_{out} = V_{in} \cdot \frac{R_2}{R_1 + R_2}
\]

\textbf{Objetivo:} $V_{in} = 5\,\text{V}$, $V_{out} \leq 3.3\,\text{V}$

\textbf{Valores seleccionados:} $R_1 = 1\,\text{k}\Omega$, $R_2 = 2\,\text{k}\Omega$

\textbf{Verificación:}
\[
V_{out} = 5 \cdot \frac{2000}{1000 + 2000} = 5 \cdot \frac{2}{3} = 3.33\,\text{V} \quad \checkmark
\]

Este valor de 3.33 V es seguro para el GPIO del Pico W (máximo 3.3 V, con un margen de tolerancia del 1\%) y suficientemente alto para ser reconocido como nivel lógico HIGH (umbral típico $V_{IH} = 2.0\,\text{V}$).

\textbf{Corriente a través del divisor:}
\[
I = \frac{V_{in}}{R_1 + R_2} = \frac{5}{3000} \approx 1.67\,\text{mA}
\]
Valor despreciable, compatible con la capacidad de corriente del pin ECHO del HC-SR04.

\subsection{Diagrama de conexiones en protoboard}

Cada divisor de tensión se implementa del siguiente modo, replicado tres veces:

\begin{enumerate}
    \item Pin \textbf{VCC} del HC-SR04 → riel de 5 V de la protoboard.
    \item Pin \textbf{GND} del HC-SR04 → riel GND de la protoboard.
    \item Pin \textbf{TRIG} del HC-SR04 → GPIO del Pico W (salida directa, sin circuito intermedio).
    \item Pin \textbf{ECHO} del HC-SR04 → extremo de \textbf{R1 (1 k$\Omega$)}.
    \item Extremo libre de R1 → nodo intermedio en protoboard. Desde este mismo nodo:
    \begin{itemize}
        \item Conectar \textbf{R2 (2 k$\Omega$)} hacia GND.
        \item Conectar un cable hacia el \textbf{pin GPIO del Pico W} (entrada ECHO\_IN).
    \end{itemize}
\end{enumerate}

\subsection{Componentes necesarios por cada divisor (×3)}

\begin{table}[H]
\centering
\renewcommand{\arraystretch}{1.3}
\begin{tabular}{|l|c|c|}
\hline
\textbf{Componente} & \textbf{Valor} & \textbf{Cantidad} \\
\hline
Resistencia & 1 k$\Omega$ / 0.25 W & 1 \\
Resistencia & 2 k$\Omega$ / 0.25 W & 1 \\
Cable jumper macho-macho & — & 4 \\
\hline
\end{tabular}
\caption{Componentes por divisor de tensión (se requieren 3 conjuntos)}
\end{table}

\subsection{Estado de avance}

\textbf{Diseño: 100\%.} Los valores están calculados, verificados y listos para implementar.\\
\textbf{Implementación física: pendiente} (sujeta a la adquisición de los sensores HC-SR04).

\section{Circuito 4 — Sistema de alimentación con regulador step-down}

\subsection{Problema que resuelve}

El sistema tiene \textbf{dos dominios de voltaje}:

\begin{itemize}
    \item \textbf{7.4 V} (batería 18650): requerido directamente por los motores DC a través del pin VM del driver TB6612FNG.
    \item \textbf{5 V regulados}: requeridos por el Pico W (pin VSYS, acepta 1.8--5.5 V), los sensores HC-SR04 (Vcc = 5 V) y el módulo GPS NEO-6M (Vcc = 3.3 V, regulado internamente desde el pin 3V3OUT del Pico W).
\end{itemize}

No es posible alimentar directamente la lógica con 7.4 V sin dañar los componentes. Se requiere un regulador step-down (convertidor buck) que reduzca eficientemente 7.4 V a 5 V.

\subsection{Módulo seleccionado}

Módulo regulador step-down genérico basado en el integrado \textbf{MP1584} o similar. Parámetros relevantes:

\begin{itemize}
    \item Voltaje de entrada: 4.5--28 V (compatible con 7.4 V de la batería).
    \item Voltaje de salida: ajustable; configurado a \textbf{5 V} mediante el potenciómetro del módulo.
    \item Corriente de salida máxima: 3 A.
    \item Eficiencia típica: 92\%.
\end{itemize}

\subsection{Procedimiento de diseño y cálculo de corriente}

Se estima la corriente total del dominio de 5 V:

\begin{table}[H]
\centering
\renewcommand{\arraystretch}{1.3}
\begin{tabular}{|l|c|}
\hline
\textbf{Componente} & \textbf{Corriente típica} \\
\hline
Raspberry Pi Pico W (WiFi activo) & 100 mA \\
HC-SR04 × 3 & 15 mA × 3 = 45 mA \\
NEO-6M GPS & 45 mA \\
\hline
\textbf{Total dominio 5 V} & \textbf{190 mA} \\
\hline
\end{tabular}
\caption{Consumo estimado del dominio de lógica}
\end{table}

El regulador seleccionado (3 A) entrega un margen de seguridad mayor a 15:1 sobre la carga de lógica. Los motores se alimentan directamente desde la batería a través del TB6612FNG, sin pasar por el regulador.

\subsection{Diagrama de conexiones en protoboard}

\begin{enumerate}
    \item \textbf{Batería (+)} → pin IN+ del módulo step-down \textbf{y} pin VM del TB6612FNG.
    \item \textbf{Batería (-)} → riel GND común de la protoboard.
    \item \textbf{Step-down OUT+} (5 V) → riel de 5 V de la protoboard.
    \item \textbf{Riel 5 V} → pin VSYS del Pico W.
    \item \textbf{Riel 5 V} → pin VCC de cada HC-SR04 (×3).
    \item \textbf{Pin 3V3OUT del Pico W} → pin VCC del módulo GPS NEO-6M.
    \item \textbf{Riel GND común} → GND de todos los componentes del sistema.
\end{enumerate}

\textbf{Nota crítica:} todos los componentes deben compartir el mismo GND (tierra común). Sin tierra común, las señales lógicas no tienen referencia y el sistema no funciona correctamente.

\subsection{Estado de avance}

\textbf{Diseño: 100\%.} El dimensionamiento está completo y el módulo step-down está incluido en el presupuesto.\\
\textbf{Implementación física: pendiente} (sujeta a la adquisición del módulo).

% ============================================================
% CAPÍTULO 4 — DIAGRAMA DE BLOQUES DEL HARDWARE
% ============================================================
\chapter{Diagrama de Bloques del Hardware}

\section{Decisión de diseño: un diagrama jerárquico en tres niveles}

El hardware del sistema SmartRoad Rover integra cinco bloques funcionales (alimentación, sensado, procesamiento, actuación y comunicación), cada uno con señales, niveles de voltaje y protocolos propios. Presentar todo el detalle en un único diagrama produciría un esquema ilegible a escala de página. Por esta razón, y siguiendo la práctica estándar en diseño de sistemas embebidos, se adopta una \textbf{jerarquía de tres diagramas}:

\begin{itemize}
    \item \textbf{Diagrama 1 — Vista general (top-level):} muestra los cinco bloques principales y sus relaciones. No incluye pines ni señales individuales. Es el mapa unificador del sistema.
    \item \textbf{Diagrama 2 — Bloque de sensado y acondicionamiento:} detalla los tres sensores HC-SR04, sus circuitos de adaptación de nivel lógico (divisores de tensión), el módulo GPS y la asignación de pines GPIO del Pico W.
    \item \textbf{Diagrama 3 — Bloque de actuación y alimentación:} detalla el driver TB6612FNG, los dos motores DC, el sistema de alimentación con sus dos dominios de voltaje y la distribución de tierra común.
\end{itemize}

Esta estructura cumple con el requerimiento del enunciado: existe un diagrama que relaciona los demás (Diagrama 1) y cada diagrama de detalle está explicado textualmente en las secciones siguientes.
Los diagramas fueron realizados con ayuda de Notebook LM.

\section{Asignación de pines GPIO del Pico W}

Antes de describir los diagramas, se fija la asignación definitiva de pines del microcontrolador. Esta tabla es la referencia para el firmware y para el montaje en protoboard:

\begin{table}[H]
\centering
\renewcommand{\arraystretch}{1.3}
\begin{tabular}{|l|c|c|l|}
\hline
\textbf{Señal} & \textbf{GPIO} & \textbf{Pin físico} & \textbf{Dirección / Nota} \\
\hline
TRIG sensor frontal  & GP0  & 1  & Salida digital \\
ECHO sensor frontal  & GP1  & 2  & Entrada — viene del divisor de tensión \\
TRIG sensor izquierdo & GP2 & 4  & Salida digital \\
ECHO sensor izquierdo & GP3 & 5  & Entrada — viene del divisor de tensión \\
TRIG sensor derecho  & GP4  & 6  & Salida digital \\
ECHO sensor derecho  & GP5  & 7  & Entrada — viene del divisor de tensión \\
GPS TX → Pico RX     & GP9  & 12 & UART1 RX \\
GPS RX → Pico TX     & GP8  & 11 & UART1 TX \\
IN1 Motor A (izq.)   & GP10 & 14 & Salida digital — dirección \\
IN2 Motor A (izq.)   & GP11 & 15 & Salida digital — dirección \\
PWM Motor A (izq.)   & GP12 & 16 & Salida PWM — velocidad \\
IN1 Motor B (der.)   & GP13 & 17 & Salida digital — dirección \\
IN2 Motor B (der.)   & GP14 & 19 & Salida digital — dirección \\
PWM Motor B (der.)   & GP15 & 20 & Salida PWM — velocidad \\
STBY TB6612FNG       & GP16 & 21 & Salida digital — enable global \\
\hline
\end{tabular}
\caption{Asignación definitiva de pines GPIO del Raspberry Pi Pico W}
\end{table}

\section{Diagrama 1 — Vista general del sistema}

\begin{figure}[H]
    \centering
    \includegraphics[width=0.90\linewidth]{diagrama1_toplevel.png}
    \caption{Diagrama de bloques top-level del sistema SmartRoad Rover}
\end{figure}

\textbf{Descripción:} el sistema se organiza alrededor del bloque de procesamiento central (Raspberry Pi Pico W). Este recibe señales del bloque de sensado (sensores ultrasónicos y GPS), genera señales de control hacia el bloque de actuación (driver de motores y motores DC), se comunica bidireccionalmente con el usuario a través del bloque de comunicación WiFi, y recibe energía del bloque de alimentación.

Las relaciones entre bloques son:
\begin{itemize}
    \item \textbf{Alimentación → Procesamiento:} suministra 5 V regulados al pin VSYS del Pico W.
    \item \textbf{Alimentación → Actuación:} suministra 7.4 V directos al pin VM del driver TB6612FNG.
    \item \textbf{Sensado → Procesamiento:} entrega señales digitales de distancia (ECHO adaptado) y tramas NMEA del GPS vía UART.
    \item \textbf{Procesamiento → Actuación:} genera señales PWM de velocidad y señales digitales de dirección hacia el driver.
    \item \textbf{Procesamiento - Comunicación:} intercambio bidireccional vía WiFi 2.4 GHz para recepción de comandos del usuario y envío de estado del sistema.
\end{itemize}

\section{Diagrama 2 — Bloque de sensado y acondicionamiento}

\begin{figure}[H]
    \centering
    \includegraphics[width=0.70\linewidth]{diagrama2_sensado.png}
    \caption{Diagrama de bloques del bloque de sensado — HC-SR04 ×3, divisores de tensión y módulo GPS}
\end{figure}

\textbf{Descripción:} cada sensor HC-SR04 genera dos señales hacia el Pico W. La señal TRIG es una salida digital del microcontrolador de 3.3 V, compatible directamente con el umbral de entrada del sensor, por lo que va sin circuito intermedio. La señal ECHO regresa del sensor a 5 V lógicos y debe ser adaptada mediante el divisor de tensión (R1 = 1 k$\Omega$, R2 = 2 k$\Omega$) antes de llegar al GPIO de entrada del Pico W, garantizando 3.33 V en el nodo intermedio.

El módulo GPS NEO-6M se comunica por UART a 3.3 V, compatible directamente con el Pico W. Su pin TX se conecta al GP9 (UART1 RX del Pico) y su pin RX al GP8 (UART1 TX). El GPS se alimenta desde el pin 3V3OUT del Pico W (3.3 V regulados internos).

Los tres sensores ultrasónicos y el GPS comparten el riel GND común de la protoboard.

\section{Diagrama 3 — Bloque de actuación y sistema de alimentación}

\begin{figure}[H]
    \centering
    \includegraphics[width=0.90\linewidth]{diagrama3_actuacion.png}
    \caption{Diagrama de bloques del bloque de actuación y sistema de alimentación}
\end{figure}

\textbf{Descripción:} el driver TB6612FNG recibe siete señales de control del Pico W: AIN1, AIN2 y PWMA para el motor izquierdo (Motor A); BIN1, BIN2 y PWMB para el motor derecho (Motor B); y STBY para habilitar o deshabilitar globalmente el driver. Las señales INx determinan la dirección de giro (avance, retroceso, frenado) y las señales PWM determinan la velocidad mediante ciclo de trabajo variable.

El driver recibe su alimentación de potencia (pin VM) directamente desde la batería de 7.4 V, sin pasar por el regulador step-down. Su alimentación lógica (pin VCC) viene del pin 3V3OUT del Pico W (3.3 V).

El sistema de alimentación distribuye la energía así: la batería alimenta en paralelo el pin VM del driver (7.4 V, dominio de potencia) y el módulo step-down (que regula a 5 V para el dominio de lógica). Desde el step-down se alimentan el Pico W (VSYS) y los tres sensores HC-SR04 (VCC). El GPS se alimenta desde el regulador interno del Pico W (3V3OUT).

Todos los componentes comparten un GND común, condición indispensable para el correcto funcionamiento de las señales digitales entre bloques.

% ============================================================
% CAPÍTULO 5 — ARQUITECTURA DE SOFTWARE
% ============================================================
\chapter{Arquitectura de Software}
 
\section{Visión general y decisiones de diseño}
 
El firmware del SmartRoad Rover se implementa en \textbf{C utilizando el Pico SDK oficial} de Raspberry Pi. La arquitectura adopta el paradigma \textbf{Polling + IRQ} exigido por el curso, donde las operaciones de tiempo crítico se manejan mediante interrupciones de hardware y la lógica de navegación y gestión de modos opera en el loop principal mediante polling. No se usa un sistema operativo en tiempo real (RTOS); en su lugar, el diseño explota correctamente los periféricos de hardware del RP2040 para lograr comportamiento determinista sin bloquear el CPU.
 
El código está organizado en \textbf{módulos funcionales independientes}, cada uno con su propio archivo fuente \texttt{.c} y su cabecera \texttt{.h}. Esta estructura facilita la depuración, el testing por módulo y el manejo del repositorio Git.
 
\section{Estructura modular del firmware}
 
\begin{table}[H]
\centering
\renewcommand{\arraystretch}{1.4}
\begin{tabular}{|l|l|p{7cm}|}
\hline
\textbf{Archivo} & \textbf{Cabecera} & \textbf{Responsabilidad} \\
\hline
\texttt{main.c} & — & Loop principal, inicialización del sistema, despacho de la FSM \\
\texttt{sensors.c} & \texttt{sensors.h} & Disparo de TRIG, ISR de ECHO, cálculo de distancia, filtrado \\
\texttt{motors.c} & \texttt{motors.h} & Control PWM de velocidad y dirección, abstracción del TB6612FNG \\
\texttt{navigation.c} & \texttt{navigation.h} & Lógica de evasión autónoma, algoritmo de selección de giro \\
\texttt{web\_server.c} & \texttt{web\_server.h} & Servidor HTTP embebido, parsing de comandos, watchdog WiFi \\
\texttt{gps.c} & \texttt{gps.h} & Recepción UART con DMA, parsing NMEA, cálculo de rumbo \\
\texttt{fsm.c} & \texttt{fsm.h} & Máquina de estados, transiciones, gestión de contingencias \\
\hline
\end{tabular}
\caption{Estructura modular del firmware}
\end{table}
 
\section{Flujo de programa: Polling + IRQ}
 
El flujo de programa sigue estrictamente el paradigma polling + interrupciones. El CPU nunca bloquea esperando eventos de hardware; en su lugar, los periféricos notifican al CPU mediante interrupciones y el CPU consulta el estado de los módulos en cada ciclo del loop.
 
\begin{itemize}
    \item \textbf{IRQ (interrupciones):} medición de tiempo de vuelo de sensores HC-SR04, recepción UART del GPS (asistida por DMA), detección de timeout del watchdog WiFi.
    \item \textbf{Polling:} lectura del resultado de las mediciones de sensores (ya calculadas por la ISR), ejecución de la lógica de navegación, despacho de la FSM, servicio de peticiones HTTP del servidor web.
\end{itemize}
 
\section{Máquina de estados (FSM)}
 
El sistema implementa una \textbf{FSM} (\textit{Finite State Machine}) principal con los siguientes estados:
 
\begin{table}[H]
\centering
\renewcommand{\arraystretch}{1.4}
\begin{tabular}{|l|p{9cm}|}
\hline
\textbf{Estado} & \textbf{Descripción} \\
\hline
\texttt{STATE\_INIT} & Inicialización de periféricos, WiFi, UART, PWM, GPIO. Transita automáticamente al primer estado seleccionado por el usuario. \\
\texttt{STATE\_MANUAL} & Ejecución de comandos recibidos por el servidor web. Watchdog activo. \\
\texttt{STATE\_AUTO\_REACTIVE} & Ciclo de polling de sensores + IRQ + lógica de evasión. \\
\texttt{STATE\_AUTO\_GPS} & Navegación hacia coordenada objetivo + evasión activa. DMA GPS activo. \\
\texttt{STATE\_SAFE\_STOP} & Detención de emergencia. Motores en cero. Espera comando de recuperación. \\
\hline
\end{tabular}
\caption{Estados de la FSM principal}
\end{table}
 
\begin{figure}[H]
    \centering
    %\includegraphics[width=0.7\linewidth]{fsm_principal.png}
    \caption{Diagrama FSM principal del sistema SmartRoad Rover}
\end{figure}
 
La implementación en C usa un \texttt{enum} y un \texttt{switch-case} en el loop principal:
 
\begin{lstlisting}[style=cstyle]
typedef enum {
    STATE_INIT,
    STATE_MANUAL,
    STATE_AUTO_REACTIVE,
    STATE_AUTO_GPS,
    STATE_SAFE_STOP
} system_state_t;
 
volatile system_state_t current_state = STATE_INIT;
 
int main(void) {
    system_init();
    while (true) {
        switch (current_state) {
            case STATE_INIT:         fsm_init();          break;
            case STATE_MANUAL:       fsm_manual();        break;
            case STATE_AUTO_REACTIVE:fsm_auto_reactive(); break;
            case STATE_AUTO_GPS:     fsm_auto_gps();      break;
            case STATE_SAFE_STOP:    fsm_safe_stop();     break;
        }
    }
}
\end{lstlisting}
 
\section{Módulo de sensores — \texttt{sensors.c}}
 
\subsection{Principio de funcionamiento}
 
La medición de distancia con el HC-SR04 requiere medir el tiempo entre el flanco de subida y el flanco de bajada de la señal ECHO con precisión de microsegundos. En C con Pico SDK esto se implementa con \textbf{interrupciones GPIO en ambos flancos}:
 
\begin{lstlisting}[style=cstyle]
/* ISR compartida para los 3 sensores HC-SR04 */
void echo_irq_handler(uint gpio, uint32_t events) {
    uint32_t now = time_us_32();
    uint8_t idx = gpio_to_sensor_index(gpio);
 
    if (events & GPIO_IRQ_EDGE_RISE) {
        sensor_data[idx].t_rise = now;
    }
    if (events & GPIO_IRQ_EDGE_FALL) {
        uint32_t pulse_us = now - sensor_data[idx].t_rise;
        /* dist (cm) = tiempo_us * velocidad_sonido / 2 */
        sensor_data[idx].distance_cm = (pulse_us * 17) / 1000;
        sensor_data[idx].valid = true;
    }
}
\end{lstlisting}
 
\subsection{Filtrado de lecturas espurias (RNF-9)}
 
Para cumplir RNF-9, cada canal de sensor mantiene un buffer circular de 3 lecturas. La distancia reportada es la mediana de las tres últimas mediciones válidas. Esto descarta picos por reflexiones irregulares sin introducir latencia significativa:
 
\begin{lstlisting}[style=cstyle]
typedef struct {
    uint32_t t_rise;
    uint16_t distance_cm;
    uint16_t buffer[SENSOR_BUF_SIZE]; /* buffer circular */
    uint8_t  buf_idx;
    bool     valid;
} sensor_channel_t;
 
uint16_t sensors_get_filtered(uint8_t idx) {
    return median3(sensor_data[idx].buffer);
}
\end{lstlisting}
 
\subsection{Disparo de TRIG (polling)}
 
El disparo del pulso TRIG es síncrono y se realiza desde el loop principal mediante polling. El tiempo de espera de 10 µs se implementa con \texttt{busy\_wait\_us(10)} del SDK, que es determinista y no bloquea el sistema de interrupciones:
 
\begin{lstlisting}[style=cstyle]
void sensors_trigger_all(void) {
    const uint trig_pins[] = {GP0, GP2, GP4};
    for (int i = 0; i < NUM_SENSORS; i++) {
        gpio_put(trig_pins[i], 1);
        busy_wait_us(10);
        gpio_put(trig_pins[i], 0);
    }
}
\end{lstlisting}
 
\section{Módulo de motores — \texttt{motors.c}}
 
El control de velocidad usa \textbf{PWM de hardware} del RP2040 a través del SDK. Cada motor tiene asignado un slice PWM independiente. La frecuencia de PWM se configura a 20 kHz (inaudible y adecuada para motores DC pequeños):
 
\begin{lstlisting}[style=cstyle]
void motors_init(void) {
    /* Motor A: GP12 - Motor B: GP15 */
    gpio_set_function(GP12, GPIO_FUNC_PWM);
    gpio_set_function(GP15, GPIO_FUNC_PWM);
    uint slice_a = pwm_gpio_to_slice_num(GP12);
    uint slice_b = pwm_gpio_to_slice_num(GP15);
    pwm_config cfg = pwm_get_default_config();
    pwm_config_set_clkdiv(&cfg, 6.25f);   /* 20 kHz */
    pwm_config_set_wrap(&cfg, 999);        /* resolucion 0-999 */
    pwm_init(slice_a, &cfg, true);
    pwm_init(slice_b, &cfg, true);
}
 
void motors_set(motor_dir_t dir_a, uint16_t speed_a,
                motor_dir_t dir_b, uint16_t speed_b) {
    gpio_put(GP10, (dir_a == FORWARD) ? 1 : 0);
    gpio_put(GP11, (dir_a == FORWARD) ? 0 : 1);
    gpio_put(GP13, (dir_b == FORWARD) ? 1 : 0);
    gpio_put(GP14, (dir_b == FORWARD) ? 0 : 1);
    pwm_set_gpio_level(GP12, speed_a);
    pwm_set_gpio_level(GP15, speed_b);
}
\end{lstlisting}
 
\section{Módulo de navegación — \texttt{navigation.c}}
 
Implementa el algoritmo de evasión autónoma. Se ejecuta en el loop principal (polling) consumiendo los datos ya calculados por las ISR de sensores. No bloquea:
 
\begin{lstlisting}[style=cstyle]
void navigation_step(void) {
    uint16_t front = sensors_get_filtered(FRONT);
    uint16_t left  = sensors_get_filtered(LEFT);
    uint16_t right = sensors_get_filtered(RIGHT);
 
    if (front > OBSTACLE_THRESHOLD_CM) {
        motors_set(FORWARD, SPEED_FULL, FORWARD, SPEED_FULL);
    } else {
        motors_set(STOP, 0, STOP, 0);
        busy_wait_ms(50);          /* pausa breve, no bloqueante para IRQ */
        if (left >= right) {
            motors_set(REVERSE, SPEED_TURN, FORWARD, SPEED_TURN);
        } else {
            motors_set(FORWARD, SPEED_TURN, REVERSE, SPEED_TURN);
        }
        busy_wait_ms(300);
    }
}
\end{lstlisting}
 
\begin{figure}[H]
    \centering
    %\includegraphics[width=0.80\linewidth]{flowchart_auto_reactive.png}
    \caption{Diagrama de flujo — modo autónomo reactivo (Polling + IRQ)}
\end{figure}
 
\section{Módulo GPS con DMA — \texttt{gps.c}}
 
\subsection{Uso de DMA (reporte explícito)}
 
El módulo GPS usa \textbf{DMA} para la recepción de tramas NMEA por UART1. El DMA transfiere bytes del registro de datos de UART1 a un buffer circular en RAM sin intervención del CPU. Cuando el DMA completa un bloque (o cuando la ISR de UART detecta el carácter \texttt{`\textbackslash n'} que termina una trama NMEA), se activa una interrupción de DMA que señaliza al loop principal que hay una trama lista para procesar:
 
\begin{lstlisting}[style=cstyle]
void gps_init(void) {
    uart_init(UART_GPS, GPS_BAUD);
    gpio_set_function(GP8, GPIO_FUNC_UART); /* TX */
    gpio_set_function(GP9, GPIO_FUNC_UART); /* RX */
 
    dma_channel = dma_claim_unused_channel(true);
    dma_channel_config cfg = dma_channel_get_default_config(dma_channel);
    channel_config_set_transfer_data_size(&cfg, DMA_SIZE_8);
    channel_config_set_read_increment(&cfg, false);
    channel_config_set_write_increment(&cfg, true);
    channel_config_set_dreq(&cfg, DREQ_UART1_RX);
 
    dma_channel_configure(dma_channel, &cfg,
        gps_buffer,              /* destino: buffer en RAM */
        &uart_get_hw(UART_GPS)->dr, /* fuente: registro UART */
        GPS_BUF_SIZE,
        true);
    dma_channel_set_irq0_enabled(dma_channel, true);
    irq_set_exclusive_handler(DMA_IRQ_0, gps_dma_irq_handler);
    irq_set_enabled(DMA_IRQ_0, true);
}
\end{lstlisting}
 
\subsection{Watchdog de señal GPS (RNF-6)}
 
Un timer de hardware del RP2040 registra el timestamp de la última trama NMEA válida recibida. El loop principal verifica en cada ciclo si han transcurrido más de 3 segundos sin trama válida y, de ser así, activa la transición de contingencia a \texttt{STATE\_SAFE\_STOP} seguido de \texttt{STATE\_AUTO\_REACTIVE}:
 
\begin{lstlisting}[style=cstyle]
bool gps_signal_valid(void) {
    return (time_us_64() - gps_last_valid_us) < GPS_TIMEOUT_US;
}
\end{lstlisting}
 
\section{Módulo servidor web — \texttt{web\_server.c}}
 
Implementa un servidor HTTP embebido sobre el stack lwIP del Pico W SDK. Sirve una página HTML con controles de movimiento. El parsing de comandos es no bloqueante: usa el modelo de callbacks de lwIP, que notifica al loop principal cuando hay datos disponibles, sin que el CPU espere activamente.
 
\subsection{Watchdog WiFi (RNF-5)}
 
Un timer de hardware cuenta el tiempo desde el último comando HTTP recibido. Si supera 500 ms, el sistema ejecuta \texttt{motors\_stop()} y transita a \texttt{STATE\_SAFE\_STOP}:
 
\begin{lstlisting}[style=cstyle]
static repeating_timer_t wifi_watchdog_timer;
 
bool wifi_watchdog_callback(repeating_timer_t *t) {
    uint64_t now = time_us_64();
    if ((now - last_cmd_timestamp_us) > WIFI_TIMEOUT_US) {
        motors_set(STOP, 0, STOP, 0);
        current_state = STATE_SAFE_STOP;
    }
    return true; /* mantener el timer activo */
}
\end{lstlisting}
 
\section{Identificación de interrupciones del sistema}
 
\begin{table}[H]
\centering
\renewcommand{\arraystretch}{1.4}
\begin{tabular}{|l|l|p{6cm}|}
\hline
\textbf{Fuente IRQ} & \textbf{Tipo} & \textbf{Acción en el handler} \\
\hline
GPIO GP1 (ECHO frontal) & Flanco ↑ y ↓ & Captura \texttt{time\_us\_32()}, calcula distancia \\
GPIO GP3 (ECHO izq.) & Flanco ↑ y ↓ & Ídem \\
GPIO GP5 (ECHO der.) & Flanco ↑ y ↓ & Ídem \\
DMA\_IRQ\_0 & Transferencia completa & Señaliza trama GPS disponible, recarga DMA \\
Repeating timer (WiFi WD) & Periódica 100 ms & Verifica timeout de último comando HTTP \\
Repeating timer (GPS WD) & Periódica 1 s & Verifica timeout de señal GPS \\
\hline
\end{tabular}
\caption{Inventario de interrupciones del firmware}
\end{table}
 
\section{Reporte de uso de DMA y máquinas de estado}
 
\begin{itemize}
    \item \textbf{DMA:} se usa \textbf{un canal DMA} para la recepción de tramas NMEA del módulo GPS por UART1. Esto libera completamente al CPU de la recepción byte a byte, reduciendo la carga de procesamiento en el dominio GPS y garantizando que ninguna trama se pierda aunque el CPU esté ejecutando lógica de navegación.
    \item \textbf{Máquina de estados:} se implementa \textbf{una FSM principal} con 5 estados (\texttt{STATE\_INIT}, \texttt{STATE\_MANUAL}, \texttt{STATE\_AUTO\_REACTIVE}, \texttt{STATE\_AUTO\_GPS}, \texttt{STATE\_SAFE\_STOP}), gestionada en \texttt{fsm.c}. La variable de estado es \texttt{volatile} para garantizar visibilidad entre el contexto de interrupción y el loop principal.
\end{itemize}
 
\section{Librerías del Pico SDK utilizadas}
 
\begin{table}[H]
\centering
\renewcommand{\arraystretch}{1.4}
\begin{tabular}{|l|p{9cm}|}
\hline
\textbf{Librería SDK} & \textbf{Uso en el proyecto} \\
\hline
\texttt{pico/stdlib.h} & Funciones base, delays, stdio \\
\texttt{hardware/gpio.h} & Configuración de pines, IRQ de GPIO \\
\texttt{hardware/pwm.h} & Control PWM de motores (GP12, GP15) \\
\texttt{hardware/uart.h} & Comunicación UART con módulo GPS \\
\texttt{hardware/dma.h} & Canal DMA para recepción GPS \\
\texttt{hardware/timer.h} & \texttt{time\_us\_32()}, \texttt{time\_us\_64()}, repeating timers \\
\texttt{pico/cyw43\_arch.h} & WiFi del Pico W (chip CYW43439) \\
\texttt{lwip/tcp.h} & Stack TCP/IP para servidor HTTP embebido \\
\hline
\end{tabular}
\caption{Librerías del Pico SDK empleadas en el firmware}
\end{table}
% ============================================================
% CAPÍTULO 6 — MANEJO DEL REPOSITORIO
% ============================================================
\chapter{Manejo del Repositorio}

% Contenido por agregar en la siguiente pregunta

% ============================================================
% REFERENCIAS
% ============================================================
\begin{thebibliography}{9}

\bibitem{picosdk}
Raspberry Pi Foundation. \textit{Raspberry Pi Pico C/C++ SDK}. \url{https://datasheets.raspberrypi.com/pico/raspberry-pi-pico-c-sdk.pdf}

\bibitem{picow}
Raspberry Pi Foundation. \textit{Raspberry Pi Pico W Datasheet}.

\bibitem{hcsr04}
\textit{HC-SR04 Ultrasonic Sensor Datasheet}.

\bibitem{tb6612}
Toshiba. \textit{TB6612FNG Motor Driver Datasheet}.

\bibitem{neo6m}
u-blox. \textit{NEO-6M GPS Module Datasheet}.

\bibitem{barr}
Michael Barr. \textit{Embedded Systems Engineering}. O'Reilly.

\end{thebibliography}

\end{document}
